\documentclass[10pt,a4paper]{amsart}
\usepackage[margin=19mm]{geometry}
\usepackage{amsmath,amssymb,mathtools}
\usepackage[scr=rsfs,cal=euler]{mathalfa}
\usepackage{xfrac}
\usepackage[unicode,hidelinks,bookmarks=false]{hyperref}
\newcommand{\ie}{i.e.\ }
\newcommand{\cf}{cf.\ }
\newcommand{\resp}{respectively\ }
\newcommand{\Subsection}{\subsection}
\providecommand{\nobreakdash}{\nobreak\hbox{-}}
\setlength{\emergencystretch}{2em}
\multlinegap0pt
\usepackage{amssymb}
\usepackage{enumerate}
\usepackage{xcolor}
\usepackage{dsfont}
\newtheorem{lemma}{Lemma}
\newtheorem{proposition}{Proposition}
\newcommand{\vol}{{\rm Vol}}
\title{\bf Variance lower bounds for geometric functionals of rotationally invariant log-concave random polytopes}
\author{Minas Pafis, Christos Pandis}
\date{Source: arXiv:2608.21663v2 (CC BY 4.0)}
\begin{document}
\maketitle

\noindent\emph{Source and licence.} \bf Variance lower bounds for geometric functionals of rotationally invariant log-concave random polytopes. Version arXiv:2608.21663v2 (CC BY 4.0), distributed by arXiv under the Creative Commons Attribution 4.0 International licence, http://creativecommons.org/licenses/by/4.0/, as recorded in the arXiv licence metadata for this exact version and shown on the abstract page https://arxiv.org/abs/2608.21663v2. Full licence notice: \texttt{licenses/arxiv-2608.21663-CC-BY-4.0.txt}.\par\smallskip

\noindent\textit{Editorial scope.} Counted unit: the article's proposition prop:q-shell-comparison (source0.tex lines 580--585). The article's complete original proof of it is delivered with it (source0.tex lines 587--789, one \texttt{proof} environment of 203 lines). No mathematics is written, repaired or re-derived: the statement and the proof are the article's own lines, in the article's own notation, and no retained line is substituted or re-flowed.  Retained uncounted as the source-local context the proof needs: lines 558--560.  Nothing was omitted from any retained span, so the package carries no omission record.  No retained line contains a citation, so the wrapper carries no bibliography and no bibliography entry.  No retained line contains a figure environment, an image-inclusion command or drawing code, so no image is delivered and the package declares no asset dependency.  Editorial formatting only, in addition: an AMS \texttt{amsart} preamble that restates the article's own declarations and macros used by the retained lines, namely the environments \texttt{lemma}, \texttt{proposition} on separate counters, together with the standard packages those lines need.  The retained environments print in this document as Lemma 1, Proposition 1 in order of appearance, whereas the article prints the same items with the numbering of its own full document.  The complete native source of the article as distributed in the arXiv e-print for this exact version is bundled unchanged as \texttt{sources/208266/source0.tex}.
\begin{lemma}\label{lem:relations}
We have that \[\frac{h_n}{a_n}, \ \frac{a_n}{r_n} \lesssim \frac{1}{\sqrt{\log n}}, \quad \frac{h_n}{r_n} \lesssim \frac{1}{\log n}.\]  
\end{lemma}

\begin{proposition}\label{prop:q-shell-comparison}
We have that
\[
q_\mu(r_n)\asymp f_\mu(r_n)a_n^{d-1}h_n.
\]
\end{proposition}

\begin{proof}
Since \(t\mapsto \rho_t\) is concave, for every \(k\ge 0\),
\begin{equation}\label{eq:rho-concavity-step}
\rho_{t_n+k+1}-r_n
\le (k+1)(\rho_{t_n+1}-\rho_{t_n})
=(k+1)h_n.
\end{equation}
We recall the standard cap-volume estimate. If
\[
\Gamma(R,\Delta):=\{x\in\mathbb R^d:\ \|x\|_2\le R,\ x_1\ge R-\Delta\},
\quad
0<\Delta\le \frac R2,
\]
is a cap of height \(\Delta\) on the ball of radius \(R\), then
\begin{equation}\label{eq:cap-volume}
\vol_d(\Gamma(R,\Delta))
\asymp
(R\Delta)^{\frac{d-1}{2}}\Delta.
\end{equation}
Indeed,
\begin{align*}
\vol_d(\Gamma(R,\Delta))
&=
\int_{R-\Delta}^R
\vol_{d-1}\bigl(\Gamma(R,\Delta)\cap\{x_1=s\}\bigr)\,ds  \\
&=
\omega_{d-1}\int_{R-\Delta}^R
(R^2-s^2)^{\frac{d-1}{2}}\,ds.
\end{align*}
For \(s\in[R-\Delta,R]\),
\[
R^2-s^2=(R-s)(R+s)\asymp R(R-s),
\]
which gives \eqref{eq:cap-volume}.

\medskip

\noindent\textbf{Lower bound.}
Let
\[
A_{n,0}:=
\{x:\langle x,e_1\rangle\ge r_n\}
\cap
\bigl(R_{t_n+1}(\mu)\setminus R_{t_n}(\mu)\bigr).
\]
Since
\[
R_{t_n}(\mu)=r_nB_2^d,
\quad
R_{t_n+1}(\mu)=(r_n+h_n)B_2^d,
\]
we have
\[
\vol_d(A_{n,0})
=
\vol_d(\Gamma(r_n+h_n,h_n)).
\]
Because \(h_n/r_n\to0\) by Lemma~\ref{lem:relations}, \eqref{eq:cap-volume} yields
\begin{equation}\label{eq:vol-lower-shell}
\vol_d(A_{n,0})
\asymp
(r_nh_n)^{\frac{d-1}{2}}h_n
=
a_n^{d-1}h_n.
\end{equation}
Moreover, for every \(x\in A_{n,0}\),
\[
e^{-1}f_\mu(r_n)\le f_\mu(x)\le f_\mu(r_n),
\]
because \(x\in R_{t_n+1}(\mu)\setminus R_{t_n}(\mu)\). Hence
\[
\mu(A_{n,0})
\asymp
f_\mu(r_n)\vol_d(A_{n,0})
\asymp
f_\mu(r_n)a_n^{d-1}h_n.
\]
Since \(A_{n,0}\subseteq \{x:\langle x,e_1\rangle\ge r_n\}\), we obtain
\begin{equation}\label{eq:q-lower}
q_\mu(r_n)\gtrsim f_\mu(r_n)a_n^{d-1}h_n.
\end{equation}


\noindent\textbf{Upper bound.}
For \(k\ge0\), define
\[
A_{n,k}:=
\{x:\langle x,e_1\rangle\ge r_n\}
\cap
\bigl(R_{t_n+k+1}(\mu)\setminus R_{t_n+k}(\mu)\bigr).
\]
Then
\[
\{x:\langle x,e_1\rangle\ge r_n\}
=
\bigcup_{k=0}^\infty A_{n,k}
\]
as a disjoint union and therefore
\[
q_\mu(r_n)=\sum_{k=0}^\infty \mu(A_{n,k}).
\]
For \(x\in A_{n,k}\), we have
\begin{equation}\label{eq:density-upper-k}
f_\mu(x)\le e^{-k}f_\mu(r_n).
\end{equation}
Also, by \eqref{eq:rho-concavity-step},
\begin{equation}\label{eq:rho-upper-k}
\rho_{t_n+k+1}-r_n\le (k+1)h_n.
\end{equation}
Hence
\[
A_{n,k}
\subseteq
\Gamma\bigl(\rho_{t_n+k+1},\,\rho_{t_n+k+1}-r_n\bigr).
\]
Let
\[
K_n:=\left\lfloor \frac{r_n}{4h_n}\right\rfloor.
\]
If \(0\le k\le K_n-1\), then, 
\[
\rho_{t_n+k+1}-r_n\le (k+1)h_n\le \frac{r_n}{4}.
\]
Therefore \eqref{eq:cap-volume} gives
\[
\vol_d(A_{n,k})
\lesssim
\bigl(r_n(k+1)h_n\bigr)^{\frac{d-1}{2}}(k+1)h_n
=
(k+1)^{\frac{d+1}{2}}a_n^{d-1}h_n.
\]
Together with \eqref{eq:density-upper-k},
\begin{equation}\label{eq:upper-small-k}
\mu(A_{n,k})
\lesssim
e^{-k}(k+1)^{\frac{d+1}{2}}
f_\mu(r_n)a_n^{d-1}h_n.
\end{equation}
Thus
\[
\sum_{k=0}^{K_n-1}\mu(A_{n,k})
\lesssim
f_\mu(r_n)a_n^{d-1}h_n,
\]
because
\[
\sum_{k=0}^\infty e^{-k}(k+1)^{\frac{d+1}{2}}<+\infty.
\]
If \(k \geq K_n\), then
\[
A_{n,k}\subseteq R_{t_n+k+1}(\mu).
\]
By \eqref{eq:rho-upper-k},
\[
\rho_{t_n+k+1}\le r_n+(k+1)h_n\leq 5(k+1)h_n
\]
for all \(k \geq K_n\). Hence
\[
\vol_d(A_{n,k})
\lesssim
(k+1)^d h_n^d.
\]
Again by \eqref{eq:density-upper-k},
\begin{equation}\label{eq:upper-large-k}
\mu(A_{n,k})
\lesssim
e^{-k}(k+1)^d f_\mu(r_n)h_n^d.
\end{equation}
Therefore
\[
\sum_{k=K_n}^{\infty}\mu(A_{n,k})
\lesssim
f_\mu(r_n)h_n^d
\sum_{k=K_n}^{\infty} e^{-k}(k+1)^d.
\]
The tail satisfies
\[
\sum_{k=K_n}^{\infty} e^{-k}(k+1)^d
\lesssim e^{-K_n/2}.
\]
 Hence
\[
\frac{
f_\mu(r_n)h_n^d
\sum_{k=K_n}^{\infty} e^{-k}(k+1)^d
}{
f_\mu(r_n)r_n^{\frac{d-1}{2}}h_n^{\frac{d+1}{2}}
}
\lesssim
e^{-K_n/2}
\left(\frac{h_n}{r_n}\right)^{\frac{d-1}{2}} \lesssim 1,
\]  due to Lemma~\ref{lem:relations}.
Thus
\[
\sum_{k=K_n}^{\infty}\mu(A_{n,k})
\lesssim f_\mu(r_n)a_n^{d-1}h_n.
\]
Combining the estimates gives
\begin{equation}\label{eq:q-upper}
q_\mu(r_n)\lesssim f_\mu(r_n)a_n^{d-1}h_n.
\end{equation}
Together with \eqref{eq:q-lower}, this proves the claim.
\end{proof}

\end{document}
